\chapter{Reproducing the implemented investigation}
\label{app:current-interpretation}

This chapter records the interpretation and integration capabilities accepted
through round four on 24 September 2026. It complements the explanations in
the main volume with exact entry points and retained evidence. Earlier scripted
increments are historical records in Appendix~\ref{app:full-implementation-record}.
Later implementation work has its own dated reports; the counts here belong to
this checkpoint and must not be read as a live resource balance.

\section{Which route is implemented}

The service can generate isolated readings, investigate alternatives with bounded
breadth-first or UCT search, check independently collected source claims, preserve
unsupported proposals and construct draft Java controls. The source inventories
are collected separately from candidate generation; independence of requests
does not establish statistical independence of errors. Composition combines
explicitly assigned component controls. Executed scenarios can rank questions
when their facts and outcomes are comparable. Registered official examples can
supply conditional arguments. A local scheduler revisits affected controls after
source changes.

The evidence addresses three targets: retained English and applicability,
proposed interpretation and formal expression, and generated Java behaviour.
Exact bytes and quotes support the first target's identity. Model judgments
and source comparisons inform the second. Java replay checks the third within
the exercised domain. No completion flag promotes one target into another.
The full scope statement is retained at
\path{docs/implementation/interpretation-round4/method-boundaries.md}.

Rounds two and three contain live development work. The accepted round-three
investigations retained 90 claim/candidate results for 23EC46 and 384 for 24EC16.
Of the latter, 224 are deterministic unavailable-program findings for seven
unencodable proposals. They are not model judgments about legal entailment.
The other 160 received source comparisons. All 474 pairs were accounted for,
with both investigations unresolved and ineligible for release. Additional
concerns and residual questions include programming, source and factual gaps;
their counts are not counts of independent legal ambiguities.

Round four revalidated this evidence and used controlled or retained inputs,
with no new model calls. Its final regression record contains 503 tests, zero
failures, errors or skips, and two existing dependency warnings. The source-update
exercise also ran a separate Java host. Suite counts from successive phases
overlap; summing them would misdescribe the evidence.

\section{Resolve an authority edition}

The catalog at \path{.localresources/sfc-authorities/b1/catalog.json} links an
issuing authority to editions, provision locations and supported aliases. A
citation in a candidate must resolve to a retained edition and exact text, not
merely a similar title. The catalog records publication and date evidence
separately from the requested assessment time. The accepted exercise covered
four official documents and six locator/date resolutions, including the
September 2023 and January 2026 Code editions.

The FAQ's title and its later answer updates require special care. The retained
page is not an independently archived 2010 edition. Its historical applicability
remains conditional, and remote statutory definitions remain unacquired.
The extraction checks compare two text extractions and preserve disagreement
or image warnings. OCR is not implemented. Matching extraction results cannot
establish that a scanned qualification has been read.

An assurance source manifest may name \texttt{authority\_registry},
\texttt{example\_registry} and \texttt{example\_scope}; relative paths resolve
against that manifest. The scope states issuer, jurisdiction and issue IDs.
The public example
\path{examples/interpretation-assurance/authority-and-examples.json}
uses the retained catalog and example registry. It is a development input,
not an institution's approved source register.

\section{Compose separate controls}

The component specification binds the source packet, inventoried claims and
complete final candidate set. Each component states its selected claims,
alternative candidate IDs, output meaning and assignment basis. The caller
must justify why the members are rival readings of one control, or separate
controls that belong together. The tool does not infer that legal grouping
from successful compilation.

The command's required arguments are:
\begin{lstlisting}
.venv/bin/python -m legalmath.cli assurance-compose \
  --investigation INVESTIGATION_DIRECTORY \
  --spec COMPONENT_SPECIFICATION.json \
  --jdk JDK_DIRECTORY --out NEW_OUTPUT_DIRECTORY
\end{lstlisting}
These capitalized arguments stand for operator-supplied paths. Use a fresh
output directory. A retained real-input specification is
\path{artifacts/interpretation/round4/B2a/attempt-01/composition/23ec46-spec.json};
its identities belong to the referenced historical investigation and cannot
be reused after changing the source or candidate.

The output includes bundles, generated Java, JARs, cases, conformance results
and a manifest. Alternative selections form separate packages. Within each
package, rules and facts use component namespaces so unrelated meanings cannot
collide. Each result retains its source claims, modality, scope, exceptions,
timing and upstream uncertainty. Missing claims, unsupported expressions,
unassigned candidates and combinations deferred by the work limit remain
explicitly incomplete. The package is a draft.

Seven builds were exercised: two controlled two-component selections and five
retained gift-control alternatives. True, false, unknown, conflict and
out-of-scope results were covered. These are not seven whole-circular
implementations; complete coverage of 23EC46 and 24EC16 remains unproved.

\section{Run a case that distinguishes readings}

The question procedure begins with a source-bound candidate comparison and its
factual witness. The record includes the source packet, candidate set, assessment
time, base candidate, subject, output convention, assumptions and source quotes.
Before replay, the implementation checks subject and output compatibility and
constructs the declared fact correspondence. The original candidate programs
then run in Python and Java on the corresponding snapshots.

The partition record retains accepted replays, excluded candidates and reasons,
work consumed and separated pairs. Unit-weight scores count executed pairs
with different projected results. They are scheduling observations, not legal
probabilities. Resource-deferred work cannot be labelled complete; unsupported
programs and incompatible facts cannot earn a score. Source-inventory omission
checks retain priority even if every comparable output agrees.

The controlled gift/discount case separated one pair and reached the repair
scheduler. The retained 23EC46 comparison had incompatible fact bindings;
24EC16 had no checked pair. Both reported \texttt{NO\_CHECKED\_DISTINCTIONS}.
This means no checked distinction was available, not that all possible readings
agreed. Conditional derived comparisons remain a separate facility. The
implemented partition and verifier are in
\path{src/legalmath/interpretation/assurance/questions.py}; retained execution
is indexed by \path{docs/implementation/interpretation-round4/b2b-result.md}.

\section{Use an official example without inventing authority}

The register at \path{.localresources/sfc-examples/b3/registry.json} binds each
example to the authority catalog, exact source quotation, issuer, edition,
proposed factors and declared outcome. The current issue is
\texttt{paragraph.3.11.applicability}. It concerns whether the paragraph applies,
not an unconditional prohibition after exceptions.

Two proposed branches of FAQ Question 1 produced twenty conditional analogy
and distinction moves on five candidates. They share one source family. Six
questions remained, including historical applicability and whether the selected
factors were legally material. The factorization is a developer proposal that
requires review; genuine quoted words do not make it exhaustive.

A generated example outside the register remains an unregistered hypothetical
proposal. A quotation cannot authenticate an invented court, case identity or
holding. This implementation has no court-case corpus, general judicial-ratio
extractor or procedure for determining binding priority. Source or registry
changes invalidate dependent analysis instead of silently retaining its former
applicability.

\section{Present proposals and preserve partial checks}

Present candidate analysis through \texttt{candidate-presentation.json}.
An entry with \texttt{WITHHELD\_CORRUPT\_PROSE} has a null presentation and points
to \texttt{prose-quality.json} and the original candidate. Raw proposals,
compiled programs and their source checks remain retained. An entry with
\texttt{NO\_PATTERN\_DETECTED} passes only the implemented pattern diagnostic.
It has no certificate of natural language coherence or legal correctness.

The detector identified five defects in the actual retained 24EC16 proposal:
placeholder subject and statement fields, self-repair chatter, and output
instructions embedded in questions. Exact authentic quotations, ordinary
uncertainty and non-English text are preserved by the checked negative cases.
The detector's coverage is deliberately narrow. Editing the proposal would
require a successor candidate and renewed checks; the presentation mechanism
does not silently rewrite its meaning.

Fidelity results are checked for exact claim/candidate pairs, source quotations
and representation quotations. The historical 23EC46 response contained 33 rows
where 32 were required; the validator identified the extra pair, and the
retained correction validated all 32. That was revalidation of an existing
correction, not a new live repair.

The per-response schema permits at most 1,000 checks and 30 concerns. The
aggregate profile bounds an investigation at 4,096 pairs, 1,080 concerns and
2,097,152 bytes; its per-call pair default is 32 with a maximum of 64. These
are engineering capacity choices. Admission considers cached and unsupported
readings as well as pending model work. It can conservatively refuse work that
might have returned shorter responses.

A completed matrix is retained in \texttt{complete.json}. Interrupted aggregation
keeps \texttt{partial.json} and the validated parts; an oversized cached part
remains separately available. Repeated concerns are retained even when a
separate presentation list deduplicates residual questions. A later failure
cannot erase earlier rows or turn a partial result into a clean interpretation.
The full profile and record names are in
\path{docs/implementation/interpretation-round4/operator-guide.md}.

\section{Freeze and execute a comparison}

The implemented adapters are \texttt{single-reader}, \texttt{isolated-readers},
\texttt{search} and \texttt{assurance}. Freeze the source, assessment time,
case-family assignments, reference outcomes and resource ceilings through
\path{legalmath.interpretation.assurance.evaluation.freeze} before dispatch.
The hidden scoring references are unavailable to the methods. Missing answers
and abstention are reported alongside unsafe and useful accepted answers.
The unit of uncertainty is a paired source-family average, as derived in
Section~\ref{sec:family-comparison} of the monograph.

A replay invocation has this form:
\begin{lstlisting}
.venv/bin/python -m legalmath.cli assurance-evaluate \
  --frozen FROZEN_STUDY.json --out NEW_OUTPUT_DIRECTORY \
  --jdk JDK_DIRECTORY --replay-responses RESPONSES.json
\end{lstlisting}
For authorized live work, the provider arguments are instead
\texttt{--allowance LEDGER.json --total-calls AUTHORIZED\_CEILING}.
The command checks the complete reservation before dispatch. Supplying the
flag does not increase the authorized ledger ceiling.

The round-four fixture reserved eight calls per method for two cases, or 64
calls. At that checkpoint only 22 of the original 100 remained. Consequently
no live comparison was run. The eight scripted jobs tested actual adapters,
hidden finite-scenario scoring and adverse outcomes. Their one development
family supplied no informative statistical separation and no estimate of
English error or reviewer savings. Additional calls alone would not make
those cases a held-out sample.

Subsequent work on references and source families is recorded separately in
\path{docs/implementation/interpretation-round5/execution-report.md}. Those reference preparations and new-source attempts form a separate
continuation. Historical call balances above must not be copied into a new run; use
the actual ledger and that run's reviewed study. A complete effectiveness study
still needs defensible outcome evidence and a justified sampling and precision
design. Internal qualified reviewers or explicit official outcomes can support
suitable tasks; there is no automatic external-consultancy requirement.

\section{Package Java and revisit source changes}

The package builder retains the draft rule bundle, policy source, JAR, host
caller, source commitments, conformance cases and manifest. A separate Java
process calls the generated public API and receives the full typed result.
Unknown or conflict must not be reduced to a Boolean authorization. The local
exercise uses draft mode and supplies no institution release identity.

The scheduler wraps the fixed assurance monitor with bounded visits and process
timeouts. Its command accepts:
\begin{lstlisting}
.venv/bin/python -m legalmath.cli assurance-schedule \
  --config REGISTRY.json --resource-root RESOURCE_DIRECTORY \
  --out NEW_OUTPUT_DIRECTORY --jdk JDK_DIRECTORY \
  --at ASSESSMENT_TIME --settings SETTINGS.json \
  --replay-responses RESPONSES.json \
  --ticks 1 --interval-seconds 60 --tick-timeout-seconds 600
\end{lstlisting}
The numbers shown illustrate a bounded replay invocation, not a recommended
bank cadence. The default is one visit. The assessment time stays fixed during
an invocation; the institution scheduler must supply the intended assessment
time for subsequent invocations. Authorized live provider arguments use the
same allowance and total-call pair as evaluation.

Before dispatch, the wrapper copies source, registry, settings and replay bytes
into the output record and binds their identities. Editing an input file while
a visit runs cannot change its frozen inputs. Failed or uncertain work consumes
an attempt and remains visible. Nested failure or unresolved meaning produces
an attention state even when the outer process exits successfully.

The final integration repair prevents an unperformed example query, a
resource-deferred partition or limited argument analysis from being cached as
fully executed. Completed but unresolved interpretation remains unresolved on
unchanged visits. Retries are bounded, and corrupted cached evidence is rejected.
The status distinguishes what ran from what was established.

The accepted update exercise launched three separate Java host processes.
Adding a fee-discount exception changed the selected restriction from true to
false on the controlled facts. An unaffected control and the earlier package
were preserved. Corrupted packages, uncertain callbacks, missed cadence and
retry exhaustion were also exercised. No actual regulatory amendment was
independently interpreted in this demonstration.

For institution staging, the bank must supply its host and JDK, service identity,
control register and fact bindings, operations-owned scheduler, update cadence
and release roles. The local command installs no service and activates no bank
policy. Confidentiality, authenticated release, audit integrity and concurrent
host transactions require evidence from that environment.

\clearpage
\section{Find the retained evidence}

The acceptance state is in
\path{artifacts/interpretation/round4/state.json}. Each accepted phase points to
an immutable attempt directory and its \texttt{run-manifest.json}; inspect the
manifest for exact commands, inputs, logs and output identities. The compact
index below explains what each record establishes.

\begin{description}[style=nextline]
\item[B0: proposal and matrix integrity]
Actual corrupt prose and the extra historical response pair, plus capacity and
partial-result retention. See \path{docs/implementation/interpretation-round4/execution-report.md}.
\item[B1: source editions]
Four documents and six locator/date resolutions, with missing definitions and
historical uncertainty retained. See \path{docs/implementation/interpretation-round4/b1-result.md}.
\item[B2a and B2b: composition and distinguishing cases]
Seven Java builds and the controlled executed partition, with no comparable
witness claimed for either retained circular. See the corresponding result
notes in \path{docs/implementation/interpretation-round4/}.
\item[B3 and B4: examples and evaluation]
One official-example family and eight scripted comparison jobs. These establish
conditional argument construction and evaluation mechanics. Their result notes
state the remaining semantic and empirical limits.
\item[B5 and B6: integration and completion]
Local Java-host update and final status/input-retention repairs. The final suite
is \path{artifacts/interpretation/round4/B6/attempt-01/tests.xml}; its adjacent
manifest binds the 503-test acceptance and operations exercise.
\end{description}

The preceding live development is documented in the round-two and round-three
execution reports under \path{docs/implementation/}. These records explain how
proposals and source judgments were obtained; round-four replay does not turn
them into fresh independent observations. The two retained circulars remain
unresolved. No method ranking, population interpretation accuracy or reduction
in review cost is established by this evidence index.
